\section{Tổng quan}
\subsection{Động lực nghiên cứu}

Sự phát triển nhanh chóng của quá trình đô thị hóa cùng với sự gia tăng liên
tục của số lượng phương tiện giao thông đã tạo ra nhiều thách thức đối với công
tác quản lý và điều hành giao thông tại các đô thị hiện đại. Các vấn đề như ùn
tắc giao thông, tai nạn đường bộ và ô nhiễm môi trường ngày càng trở nên phổ
biến, đòi hỏi các cơ quan quản lý phải có khả năng giám sát và phân tích tình
hình giao thông một cách hiệu quả. Trong bối cảnh đó, Hệ thống Giao thông Thông
minh (Intelligent Transportation Systems -- ITS) đã trở thành một hướng nghiên
cứu quan trọng nhằm nâng cao mức độ an toàn, hiệu quả và tính bền vững của hệ
thống giao thông \cite{YANG2018143}.

Trong Hệ thống Giao thông Thông minh (ITS), phát hiện và đếm phương tiện là một
trong những nhiệm vụ quan trọng, cung cấp dữ liệu phục vụ nhiều ứng dụng như
giám sát lưu lượng giao thông, điều khiển tín hiệu giao thông, phát hiện ùn
tắc, dự báo lưu lượng phương tiện và hỗ trợ quy hoạch hạ tầng giao thông
\cite{YANG2018143}. Trước đây, việc thu thập thông tin giao thông chủ yếu
dựa trên các cảm biến chuyên dụng như vòng từ cảm ứng, radar hoặc cảm biến
laser. Mặc dù các thiết bị này có khả năng cung cấp dữ liệu tương đối chính
xác, chúng thường đòi hỏi chi phí lắp đặt và bảo trì cao, đồng thời dễ bị ảnh
hưởng bởi các điều kiện môi trường \cite{YANG2018143}.

Sự phát triển và phổ biến của các hệ thống camera giám sát giao thông đã tạo ra
nguồn dữ liệu hình ảnh và video phong phú, mở ra hướng tiếp cận mới cho việc
thu thập và phân tích thông tin giao thông trên quy mô lớn. So với các loại cảm
biến truyền thống, camera có ưu thế về chi phí, tính linh hoạt và khả năng cung
cấp thông tin trực quan phục vụ nhiều nhiệm vụ phân tích khác nhau
\cite{YANG2018143}.

Nhờ sự phát triển của học sâu, các hệ thống phát hiện và đếm phương tiện dựa
trên dữ liệu video đã đạt được độ chính xác cao và có khả năng hoạt động hiệu
quả trong các bối cảnh giao thông thực tế, đặc biệt trên các tuyến đường cao
tốc và đô thị \cite{song2019}.

Vì vậy, nghiên cứu các phương pháp phát hiện và đếm phương tiện từ video giám
sát giao thông không chỉ mang ý nghĩa khoa học trong lĩnh vực thị giác máy tính
mà còn có giá trị thực tiễn đối với việc xây dựng và phát triển các hệ thống
giao thông thông minh.

\subsection{Phát biểu bài toán}

Bài toán phát hiện và đếm phương tiện trong giám sát giao thông sử dụng dữ liệu
đầu vào là video hoặc chuỗi ảnh được thu nhận từ các camera giám sát lắp đặt
tại các tuyến đường, nút giao thông hoặc đường cao tốc. Mục tiêu của hệ thống
là xác định vị trí của các phương tiện xuất hiện trong từng khung hình, theo dõi
sự di chuyển của chúng giữa các khung hình liên tiếp và thực hiện đếm số lượng
phương tiện đi qua một vùng quan sát xác định trước. Thông tin thu được từ quá
trình phát hiện và đếm phương tiện là cơ sở cho nhiều ứng dụng trong hệ thống
giao thông thông minh như giám sát lưu lượng giao thông, phát hiện ùn tắc và hỗ
trợ quản lý giao thông \cite{YANG2018143}.

Kết quả đầu ra của hệ thống bao gồm vị trí phương tiện dưới dạng hộp giới hạn,
thông tin theo dõi đối tượng và số lượng phương tiện được ghi nhận trong khoảng
thời gian quan sát. Hệ thống cần đảm bảo độ chính xác cao trong điều kiện giao
thông thực tế, đồng thời đáp ứng yêu cầu xử lý gần thời gian thực nhằm phục vụ
hiệu quả cho các ứng dụng giám sát và điều hành giao thông.
\subsection{Ý nghĩa khoa học và thực tiễn}

\subsubsection{Ý nghĩa khoa học}

Đề tài góp phần nghiên cứu khả năng ứng dụng các mô hình học sâu trong bài toán
phát hiện đối tượng và phân tích dữ liệu giao thông từ video. Việc đánh giá
hiệu quả của các mô hình phát hiện phương tiện trong những điều kiện giao thông
khác nhau giúp bổ sung cơ sở thực nghiệm cho lĩnh vực thị giác máy tính và giao
thông thông minh. Ngoài ra, đề tài cũng tạo tiền đề cho các nghiên cứu liên
quan đến theo dõi đối tượng, phân tích hành vi giao thông và dự báo lưu lượng
phương tiện trong tương lai.

\subsubsection{Ý nghĩa thực tiễn}

Hệ thống phát hiện và đếm phương tiện có thể được ứng dụng trong công tác giám
sát giao thông tự động, hỗ trợ điều khiển tín hiệu giao thông thích ứng, đánh
giá mật độ phương tiện, phát hiện ùn tắc và thu thập dữ liệu phục vụ quy hoạch
hạ tầng giao thông. Việc tự động hóa quá trình thu thập dữ liệu giao thông giúp
giảm chi phí vận hành, nâng cao độ chính xác và cung cấp thông tin theo thời
gian thực cho các cơ quan quản lý giao thông.

\subsection{Thách thức của bài toán}

Mặc dù các phương pháp học sâu đã đạt được nhiều tiến bộ trong những năm gần
đây, bài toán phát hiện và đếm phương tiện trong môi trường thực tế vẫn còn
nhiều thách thức. Điều kiện ánh sáng thay đổi giữa ban ngày và ban đêm, các
hiện tượng thời tiết bất lợi như mưa hoặc sương mù, hiện tượng che khuất giữa
các phương tiện và sự khác biệt về góc quan sát của camera đều có thể ảnh hưởng
đến độ chính xác của hệ thống. Bên cạnh đó, sự đa dạng về kích thước và chủng
loại phương tiện cũng làm tăng độ phức tạp của quá trình phát hiện và theo dõi
đối tượng \cite{YANG2018143}.

\subsection{Phạm vi nghiên cứu}

Đề tài tập trung nghiên cứu bài toán phát hiện, theo dõi và đếm phương tiện
giao thông đường bộ dựa trên dữ liệu video từ các camera giám sát cố định. Các
thí nghiệm được thực hiện trên bộ dữ liệu UA-DETRAC, bao gồm các phương tiện
như ô tô, xe buýt, xe tải nhẹ và các loại phương tiện khác được ghi nhận trong
tập dữ liệu.

Về phương pháp, đề tài sử dụng mô hình phát hiện đối tượng dựa trên học sâu để
xác định vị trí phương tiện trong từng khung hình, kết hợp với thuật toán theo
dõi đối tượng nhằm duy trì định danh của phương tiện và thực hiện đếm phương
tiện khi chúng đi qua vùng quan sát xác định trước.
